% META: {"id": "1SPE-PRODUIT-SCALAIRE-QCM", "chapitre": "1SPE-PRODUIT-SCALAIRE", "type_objet": "qcm", "genere_depuis": "chapitres/1SPE-PRODUIT-SCALAIRE/qcm/1SPE-PRODUIT-SCALAIRE-QCM.json", "status": "generated"}
% Fichier genere par scripts/build_qcm_tex.py — ne pas editer a la main.

\section*{\textcolor{chapcolor}{\MakeUppercase{Produit scalaire — QCM d'auto-évaluation}}}

\begin{center}
\textit{Pour chaque question, une seule reponse est exacte.}
\end{center}

\begin{enumerate}

\item[] \textbf{Capacite C1}

\item \textbf{[Q1]} Soient $\vec{u}(3,-2)$ et $\vec{v}(1,4)$. Que vaut $\vec{u} \cdot \vec{v}$ ?
  \begin{enumerate}[label=\Alph*.]
    \item $-5$
    \item $5$
    \item $-11$
    \item $11$
  \end{enumerate}

\item \textbf{[Q2]} Si $\|\vec{u}\|=5$, $\|\vec{v}\|=4$ et $(\vec{u},\vec{v}) = \dfrac{\pi}{3}$, quelle est la valeur de $\vec{u} \cdot \vec{v}$ ?
  \begin{enumerate}[label=\Alph*.]
    \item $20$
    \item $10$
    \item $10\sqrt{3}$
    \item $10\sqrt{2}$
  \end{enumerate}

\item \textbf{[Q3]} Le produit scalaire $\vec{u} \cdot \vec{v}$ est :
  \begin{enumerate}[label=\Alph*.]
    \item un vecteur
    \item un nombre réel
    \item toujours positif
    \item défini uniquement si $\vec{u}$ et $\vec{v}$ sont non nuls
  \end{enumerate}

\item[] \textbf{Capacite C2}

\item \textbf{[Q4]} $\Vert\vec{u}+\vec{v}\Vert^2 = ?$
  \begin{enumerate}[label=\Alph*.]
    \item $\Vert\vec{u}\Vert^2 + \Vert\vec{v}\Vert^2$
    \item $\Vert\vec{u}\Vert^2 + 2\vec{u}\cdot\vec{v} + \Vert\vec{v}\Vert^2$
    \item $\Vert\vec{u}\Vert^2 - 2\vec{u}\cdot\vec{v} + \Vert\vec{v}\Vert^2$
    \item $(\vec{u}\cdot\vec{v})^2$
  \end{enumerate}

\item \textbf{[Q5]} $\vec{u} \cdot \vec{u} = ?$
  \begin{enumerate}[label=\Alph*.]
    \item $\Vert\vec{u}\Vert$
    \item $\Vert\vec{u}\Vert^2$
    \item $2\Vert\vec{u}\Vert$
    \item $0$
  \end{enumerate}

\item \textbf{[Q6]} Sachant que $\|\vec{u}\|=3$, $\|\vec{v}\|=4$ et $\vec{u}\cdot\vec{v}=5$, quelle est la valeur de $\|\vec{u}+\vec{v}\|$ ?
  \begin{enumerate}[label=\Alph*.]
    \item $5$
    \item $\sqrt{35}$
    \item $7$
    \item $\sqrt{25}$
  \end{enumerate}

\item[] \textbf{Capacite C3}

\item \textbf{[Q7]} Deux vecteurs non nuls $\vec{u}$ et $\vec{v}$ sont orthogonaux si et seulement si :
  \begin{enumerate}[label=\Alph*.]
    \item $\vec{u}\cdot\vec{v} > 0$
    \item $\vec{u}\cdot\vec{v} < 0$
    \item $\vec{u}\cdot\vec{v} = 0$
    \item $\Vert\vec{u}\Vert = \Vert\vec{v}\Vert$
  \end{enumerate}

\item \textbf{[Q8]} Les vecteurs $\vec{u}(2,3)$ et $\vec{v}(3,-2)$ sont :
  \begin{enumerate}[label=\Alph*.]
    \item colinéaires
    \item orthogonaux
    \item de même norme
    \item de sens contraire
  \end{enumerate}

\item \textbf{[Q9]} Quel est l'angle orienté entre les vecteurs $(1,0)$ et $(1,1)$ ?
  \begin{enumerate}[label=\Alph*.]
    \item $\dfrac{\pi}{6}$
    \item $\dfrac{\pi}{4}$
    \item $\dfrac{\pi}{3}$
    \item $\dfrac{\pi}{2}$
  \end{enumerate}

\item[] \textbf{Capacite C4}

\item \textbf{[Q10]} Un point $M$ appartient à la médiatrice du segment $[AB]$ si et seulement si :
  \begin{enumerate}[label=\Alph*.]
    \item $\vec{MA} \cdot \vec{MB} = 0$
    \item $MA = MB$
    \item $\vec{MA} + \vec{MB} = \vec{0}$
    \item $MA + MB = AB$
  \end{enumerate}

\item \textbf{[Q11]} Dans un triangle de côtés $a, b$ formant un angle $\theta$, l'aire est donnée par :
  \begin{enumerate}[label=\Alph*.]
    \item $ab \cos(\theta)$
    \item $\dfrac{1}{2} ab \cos(\theta)$
    \item $\dfrac{1}{2} ab \sin(\theta)$
    \item $ab \sin(\theta)$
  \end{enumerate}

\item \textbf{[Q12]} Soit $H$ le pied de la hauteur issue de $A$ dans le triangle $ABC$. Quelle relation est exacte ?
  \begin{enumerate}[label=\Alph*.]
    \item $\vec{AH} \cdot \vec{AB} = 0$
    \item $\vec{AH} \cdot \vec{BC} = 0$
    \item $\vec{AH} \cdot \vec{AC} = 0$
    \item $AH = BH$
  \end{enumerate}

\item[] \textbf{Capacite C5}

\item \textbf{[Q13]} Quelle est l'expression exacte de la formule d'Al-Kashi ?
  \begin{enumerate}[label=\Alph*.]
    \item $a^2 = b^2 + c^2 + 2bc \cos(A)$
    \item $a^2 = b^2 + c^2 - 2bc \cos(A)$
    \item $a^2 = b^2 - c^2 - 2bc \cos(A)$
    \item $a = b + c - 2bc \cos(A)$
  \end{enumerate}

\item \textbf{[Q14]} Dans un triangle avec $b=5$, $c=8$ et $\widehat{A} = \dfrac{\pi}{3}$, que vaut $a^2$ ?
  \begin{enumerate}[label=\Alph*.]
    \item $49$
    \item $89$
    \item $129$
    \item $9$
  \end{enumerate}

\item \textbf{[Q15]} Si $\widehat{A} = \dfrac{\pi}{2}$, la formule d'Al-Kashi redonne :
  \begin{enumerate}[label=\Alph*.]
    \item $a^2 = b^2 - c^2$
    \item $a^2 = b^2 + c^2$
    \item $a^2 = 2bc$
    \item $a = b + c$
  \end{enumerate}

\end{enumerate}

\clearpage
\section*{Cle de correction — reservee au professeur}

\begin{center}
\begin{tabular}{lll}
\hline
Question & Capacite & Reponse exacte \\
\hline
Q1 & C1 & \textbf{A} \\
Q2 & C1 & \textbf{B} \\
Q3 & C1 & \textbf{B} \\
Q4 & C2 & \textbf{B} \\
Q5 & C2 & \textbf{B} \\
Q6 & C2 & \textbf{B} \\
Q7 & C3 & \textbf{C} \\
Q8 & C3 & \textbf{B} \\
Q9 & C3 & \textbf{B} \\
Q10 & C4 & \textbf{B} \\
Q11 & C4 & \textbf{C} \\
Q12 & C4 & \textbf{B} \\
Q13 & C5 & \textbf{B} \\
Q14 & C5 & \textbf{A} \\
Q15 & C5 & \textbf{B} \\
\hline
\end{tabular}
\end{center}

\subsection*{Diagnostic des reponses erronees}

\begin{itemize}
  \item \textbf{Q1}
  \begin{itemize}
    \item \textbf{B} — Oubli du signe dans $(-2) \times 4$. \emph{Renvoi : M1.}
    \item \textbf{C} — Erreur de calcul dans le produit des composantes. \emph{Renvoi : M1.}
    \item \textbf{D} — Addition des composantes au lieu du produit scalaire. \emph{Renvoi : M1.}
  \end{itemize}
  \item \textbf{Q2}
  \begin{itemize}
    \item \textbf{A} — Oubli du facteur $\cos\!\left(\dfrac{\pi}{3}\right) = \dfrac{1}{2}$. \emph{Renvoi : M1.}
    \item \textbf{C} — Confusion entre le cosinus et le sinus. \emph{Renvoi : M1; R3.}
    \item \textbf{D} — Mauvaise valeur du cosinus de l'angle. \emph{Renvoi : M1; R3.}
  \end{itemize}
  \item \textbf{Q3}
  \begin{itemize}
    \item \textbf{A} — Confusion avec le produit vectoriel. \emph{Renvoi : M1.}
    \item \textbf{C} — Le produit scalaire peut être négatif si l'angle est obtus. \emph{Renvoi : M1.}
    \item \textbf{D} — On pose $\vec{u} \cdot \vec{v} = 0$ si l'un des vecteurs est nul. \emph{Renvoi : M1.}
  \end{itemize}
  \item \textbf{Q4}
  \begin{itemize}
    \item \textbf{A} — Oubli du double produit $2\vec{u}\cdot\vec{v}$. \emph{Renvoi : M2.}
    \item \textbf{C} — Signe du double produit (c'est la formule pour $\Vert\vec{u}-\vec{v}\Vert^2$). \emph{Renvoi : M2.}
    \item \textbf{D} — Confusion entre norme du carré de la somme et carré du produit scalaire. \emph{Renvoi : M2.}
  \end{itemize}
  \item \textbf{Q5}
  \begin{itemize}
    \item \textbf{A} — Oubli du carré dans la définition du carré scalaire. \emph{Renvoi : M2.}
    \item \textbf{C} — Confusion de coefficient. \emph{Renvoi : M2.}
    \item \textbf{D} — $\vec{u}\cdot\vec{u}=0$ si et seulement si $\vec{u}=\vec{0}$. \emph{Renvoi : M2.}
  \end{itemize}
  \item \textbf{Q6}
  \begin{itemize}
    \item \textbf{A} — $\sqrt{25}=5$ mais $\|\vec{u}+\vec{v}\|^2 = 9 + 10 + 16 = 35$. \emph{Renvoi : M2.}
    \item \textbf{C} — Erreur d'application de la formule. \emph{Renvoi : M2.}
    \item \textbf{D} — $\|\vec{u}+\vec{v}\|^2 = 35$, et non $25$. \emph{Renvoi : M2.}
  \end{itemize}
  \item \textbf{Q7}
  \begin{itemize}
    \item \textbf{A} — Le signe strict ne caractérise pas l'orthogonalité. \emph{Renvoi : M3.}
    \item \textbf{B} — Idem. \emph{Renvoi : M3.}
    \item \textbf{D} — Avoir la même norme ne garantit pas l'orthogonalité. \emph{Renvoi : M3.}
  \end{itemize}
  \item \textbf{Q8}
  \begin{itemize}
    \item \textbf{A} — Le déterminant vaut $2 \times (-2) - 3 \times 3 = -13 
eq 0$. \emph{Renvoi : M3.}
    \item \textbf{C} — $\Vert\vec{u}\Vert = \Vert\vec{v}\Vert = \sqrt{13}$ est vrai mais l'orthogonalité est la propriété caractérisée. \emph{Renvoi : M3.}
    \item \textbf{D} — Option incorrecte. \emph{Renvoi : M3.}
  \end{itemize}
  \item \textbf{Q9}
  \begin{itemize}
    \item \textbf{A} — Confusion dans les valeurs remarquables de l'angle. \emph{Renvoi : M3; R3.}
    \item \textbf{C} — Confusion de valeur remarquable. \emph{Renvoi : M3; R3.}
    \item \textbf{D} — Le produit scalaire est non nul ($1 
eq 0$). \emph{Renvoi : M3.}
  \end{itemize}
  \item \textbf{Q10}
  \begin{itemize}
    \item \textbf{A} — Ce serait l'équation du cercle de diamètre $[AB]$. \emph{Renvoi : M4.}
    \item \textbf{C} — Ceci caractérise le milieu de $[AB]$. \emph{Renvoi : M4.}
    \item \textbf{D} — Ceci caractérise l'appartenance au segment $[AB]$. \emph{Renvoi : M4.}
  \end{itemize}
  \item \textbf{Q11}
  \begin{itemize}
    \item \textbf{A} — Cosinus au lieu de sinus et oubli du facteur $1/2$. \emph{Renvoi : M4.}
    \item \textbf{B} — Cosinus au lieu de sinus. \emph{Renvoi : M4.}
    \item \textbf{D} — Oubli du facteur $1/2$. \emph{Renvoi : M4.}
  \end{itemize}
  \item \textbf{Q12}
  \begin{itemize}
    \item \textbf{A} — La hauteur est perpendiculaire au côté opposé $(BC)$, pas à $(AB)$. \emph{Renvoi : M4.}
    \item \textbf{C} — La hauteur est perpendiculaire à $(BC)$, pas à $(AC)$. \emph{Renvoi : M4.}
    \item \textbf{D} — Relation de longueurs non justifiée. \emph{Renvoi : M4.}
  \end{itemize}
  \item \textbf{Q13}
  \begin{itemize}
    \item \textbf{A} — Signe $+$ au lieu de $-$ devant le terme en cosinus. \emph{Renvoi : M5.}
    \item \textbf{C} — Signe $-$ devant $c^2$ au lieu de $+$. \emph{Renvoi : M5.}
    \item \textbf{D} — Oubli des carrés sur les longueurs des côtés. \emph{Renvoi : M5.}
  \end{itemize}
  \item \textbf{Q14}
  \begin{itemize}
    \item \textbf{B} — Oubli du terme $-2bc \cos(A)$. \emph{Renvoi : M5.}
    \item \textbf{C} — Erreur de signe $+2bc \cos(A)$ au lieu de $-$. \emph{Renvoi : M5.}
    \item \textbf{D} — Erreur de calcul. \emph{Renvoi : M5.}
  \end{itemize}
  \item \textbf{Q15}
  \begin{itemize}
    \item \textbf{A} — Signe $-$ au lieu de $+$. \emph{Renvoi : M5; R5.}
    \item \textbf{C} — Formule incorrecte. \emph{Renvoi : M5; R5.}
    \item \textbf{D} — Ce n'est pas l'égalité du théorème de Pythagore. \emph{Renvoi : M5; R5.}
  \end{itemize}
\end{itemize}
